\clearpage
\begin{proofbox}
\textbf{方法一：最快证明｜过焦点代入＋韦达/判别式（考场速写）}

\textbf{信号：}题目同时出现“焦点弦长”“直线斜率 $k$”“双曲线”。
令 $F_2=(c,0)$，\textbf{只设一个未知量} $t=x-c$，于是点为 $(c+t,kt)$。
\begin{center}% a=3,b=4,c=5,k=5/12. A=(-29/7,-80/21),B=(101/33,-80/99).
\begin{tikzpicture}[font=\small,>=stealth]
\draw[gray!45,rounded corners=2pt] (0,0) rectangle (14.2,10.2);
\node at (7.1,9.55) {根差为何能求弦长：先看弦与横、竖投影组成的直角三角形};
\node at (7.1,8.65) {示例 $a=3,b=4,c=5,k=5/12$；图中 $t_A<t_B<0$};
\begin{scope}[shift={(6.1,7)},scale=.8]
\coordinate (A) at (-4.1428571429,-3.8095238095);
\coordinate (B) at (3.0606060606,-.8080808081);
\coordinate (F) at (5,0);
\coordinate (H) at (3.0606060606,-3.8095238095);
\draw[gray,->] (-5.4,0)--(6,0) node[right] {$x$};
\draw[blue!75!black,thick,domain=-4.5:1,samples=90,variable=\v] plot ({3*sqrt(1+\v*\v/16)},\v);
\draw[blue!75!black,thick,domain=-4.5:1,samples=90,variable=\v] plot ({-3*sqrt(1+\v*\v/16)},\v);
\draw[red!75!black,dashed] (B)--(F);
\draw[red!75!black,very thick] (A)--(B) node[midway,above left=4pt] {$L$};
\draw[orange!85!black,thick] (A)--(H) node[midway,below=6pt] {$|t_A-t_B|$};
\draw[purple!80!black,thick] (H)--(B);
\node[anchor=west,text=purple!80!black] at (4.8,-2.3088023088) {$|k|\,|t_A-t_B|$};
\draw[gray] (2.8606060606,-3.8095238095)--(2.8606060606,-3.6095238095)--(3.0606060606,-3.6095238095);
\fill (A) circle(2pt) node[left=10pt] {$A$};
\fill (B) circle(2pt) node[left=10pt] {$B$};
\fill (F) circle(2pt) node[above right=5pt] {$F_2$};
\fill (H) circle(2pt) node[below right=5pt] {$H$};
\end{scope}
\node at (7.1,2.25) {$\displaystyle L^2=|t_A-t_B|^2+k^2|t_A-t_B|^2$};
\node at (7.1,1.2) {$\displaystyle L=\sqrt{1+k^2}\,|t_A-t_B|$};
\node at (7.1,.5) {$t=x-c$ 是横向增量；两根都为负，只说明两个交点都在焦点左侧。};
\end{tikzpicture}
\end{center}
\textbf{先读图中的直角三角形：} 两条直角边分别是横、竖坐标差，故根差先给出横向投影。

\textbf{第1步：}代入双曲线后乘 $a^2b^2$：
\[
\boxed{(b^2-a^2k^2)t^2+2cb^2t+b^4=0}.\tag{1}
\]
\textbf{第2步：}无需解根，直接算出
\[
\Delta=(2cb^2)^2-4(b^2-a^2k^2)b^4
=\boxed{4a^2b^4(1+k^2)}.
\]
\textbf{第3步：}弦长等于“横坐标差乘斜率长度系数”：
\[
L=\sqrt{1+k^2}\,|t_1-t_2|
=\sqrt{1+k^2}\frac{\sqrt\Delta}{|b^2-a^2k^2|}
=\boxed{\frac{2ab^2(1+k^2)}{|a^2k^2-b^2|}}.
\]
前提 $b^2-a^2k^2\ne0$，也即 $k\ne\pm b/a$。
当该系数为零，(1) 退化为一次方程，只有一个有限交点。

\tcbbreak
\textbf{方法一续：先看弦长怎样变化，再比较两个下界。}
\begin{center}% Actual L(theta), a=3,b=4,c=5; horizontal degree/10, vertical .14L.
\begin{tikzpicture}[font=\small,>=stealth]
\draw[gray!45,rounded corners=2pt] (0,0) rectangle (14.2,9.2);
\node at (7.1,8.55) {转动直线看弦长：临界处无界增长，水平和竖直分别取最小值};
\node at (7.1,7.75) {图取 $a=3,b=4,c=5$，只看 $0\le\theta\le\pi/2$};
\begin{scope}[shift={(2.3,1.6)}]
\draw[gray,->] (0,0)--(9.7,0) node[right] {$\theta$};
\draw[gray,->] (0,0)--(0,5.8) node[above] {$L$};
\draw[gray,dashed] ({acos(3/5)/10},0)--({acos(3/5)/10},5.8);
\draw[red!75!black,very thick,->,domain=0:{acos(sqrt((9+96/40)/25))},samples=100,variable=\v]
 plot ({\v/10},{.14*96/(25*cos(\v)^2-9)});
\draw[purple!80!black,very thick,<-,domain={acos(sqrt((9-96/40)/25))}:90,samples=100,variable=\v]
 plot ({\v/10},{.14*96/(9-25*cos(\v)^2)});
\fill[red!75!black] (0,.84) circle(2.2pt) node[left=6pt] {$2a=6$};
\fill[purple!80!black] (9,{.14*32/3}) circle(2.2pt) node[above right=5pt] {$2b^2/a=32/3$};
\node[below=5pt] at (0,0) {$0$};
\node[below=5pt] at (9,0) {$\pi/2$};
\node[below=5pt] at ({acos(3/5)/10},0) {$\theta_0=\arccos(3/5)$};
\node[red!75!black] at (2,2.8) {异支};
\node[purple!80!black] at (7.4,3.4) {同支};
\end{scope}
\node at (7.1,.45) {临界角 $\theta_0$ 对应渐近线方向；实心点说明两个最小值都能取到。};
\end{tikzpicture}
\end{center}
图中左、右两段分别是异支、同支；两端实心点对应水平、竖直弦。
\textbf{最值两行速写：}设 $u=k^2$，$c^2=a^2+b^2$，则
\[
\boxed{\begin{aligned}
0\le u< b^2/a^2:&\quad L-2a=\frac{2ac^2u}{b^2-a^2u}\ge0,\\[3pt]
u>b^2/a^2:&\quad L-\frac{2b^2}{a}
=\frac{2b^2c^2}{a(a^2u-b^2)}>0.
\end{aligned}}
\]
第一行等号在 $k=0$；第二行无有限 $k$ 能取等，但竖直线 $x=c$ 取等。
以上既给出下界，也因趋近 $|k|=b/a$ 时 $L\to\infty$、随 $k$ 连续而覆盖所写的全部范围。
\end{proofbox}

\clearpage
\begin{proofbox}
\textbf{方法二：最低门槛｜先看分支，再看「多出来的长度」}
\begin{center}% 精确坐标：H: x^2/9-y^2/16=1, c=5, b^2/a=16/3.
% 两图中蓝色曲线由解析式 x=+-3 sqrt(1+y^2/16) 绘制，非手绘。
\begin{tikzpicture}[x=0.39cm,y=0.39cm,>=stealth, font=\small]
\begin{scope}
  \draw[->,gray] (-6.8,0)--(7.3,0) node[right]{$x$};
  \draw[->,gray] (0,-7.8)--(0,7.8) node[above]{$y$};
  \draw[thick,blue!75!black,domain=-7:7,samples=90] plot ({3*sqrt(1+\x*\x/16)},{\x});
  \draw[thick,blue!75!black,domain=-7:7,samples=90] plot ({-3*sqrt(1+\x*\x/16)},{\x});
  \draw[very thick,red!75!black] (5,-5.33333333)--(5,5.33333333);
  \fill (5,0) circle (2.1pt) node[below right]{$F_2$};
  \fill[red!75!black] (5,5.33333333) circle (1.8pt) node[above left=5pt]{$A$};
  \fill[red!75!black] (5,-5.33333333) circle (1.8pt) node[below left=5pt]{$B$};
  \node[align=center] at (0,-9.0) {同支焦点弦\quad $x=c$\\最短 $2b^2/a$};
\end{scope}
\begin{scope}[xshift=6.6cm]
  \draw[->,gray] (-6.8,0)--(7.3,0) node[right]{$x$};
  \draw[->,gray] (0,-7.8)--(0,7.8) node[above]{$y$};
  \draw[thick,blue!75!black,domain=-7:7,samples=90] plot ({3*sqrt(1+\x*\x/16)},{\x});
  \draw[thick,blue!75!black,domain=-7:7,samples=90] plot ({-3*sqrt(1+\x*\x/16)},{\x});
  \draw[very thick,red!75!black] (-3,0)--(3,0);
  \fill (5,0) circle (2.1pt) node[above right]{$F_2$};
  \fill[red!75!black] (-3,0) circle (1.8pt) node[above left=5pt]{$A$};
  \fill[red!75!black] (3,0) circle (1.8pt) node[above right=5pt]{$B$};
  \node[align=center] at (0,-9.0) {异支焦点弦\quad $y=0$\\最短 $2a$};
\end{scope}
\end{tikzpicture}
\end{center}

\textbf{为什么采用：}不用分别求 $A,B$ 坐标，不必求导，适合看到成品公式后马上判断最值。

\textbf{第1步：}记住双曲线自己的分界线斜率 $b/a$。
较平（$|k|<b/a$）是异支，较陡（$|k|>b/a$）是同支，竖直归同支。

\textbf{第2步：}异支连线必跨越 $x=-a$ 与 $x=a$：
\[
\boxed{|AB|\ge|x_A-x_B|\ge2a}.\tag{2}
\]
只有 $A=(-a,0),\ B=(a,0)$ 时等号成立。
\textbf{这一结论完全不需要弦长公式！}

\textbf{第3步：}同支用\textbf{不求导的直接比较}。设 $u=k^2>b^2/a^2$，
\[
\begin{aligned}
L&=\frac{2ab^2(1+u)}{a^2u-b^2}
=\frac{2b^2}{a}\left(1+\frac{a^2+b^2}{a^2u-b^2}\right)
>\frac{2b^2}{a}.
\end{aligned}
\]
看似严格大于，\textbf{不代表最小值取不到}：斜率公式遗漏了竖直线，
在 $x=c$ 上，$A,B=(c,\pm b^2/a)$，故 $L=2b^2/a$ 正好取等。

\textbf{极低记忆负担：}\emph{异支跨实轴，最短 $2a$；同支拉竖线，最短 $2b^2/a$。}
\textbf{务必区分：}竖直线不是 $k=\infty$ 的代数代入，而是另一个必须核验的几何情形。
\end{proofbox}

\clearpage
\begin{proofbox}
\textbf{方法三：常规解析几何｜联立二次方程＋根差＋角度统一}
\begin{center}% a=3,b=4,c=5,k=5/12. A=(-29/7,-80/21),B=(101/33,-80/99).
\begin{tikzpicture}[font=\small,>=stealth]
\draw[gray!45,rounded corners=2pt] (0,0) rectangle (14.2,10.2);
\node at (7.1,9.55) {根差为何能求弦长：先看弦与横、竖投影组成的直角三角形};
\node at (7.1,8.65) {示例 $a=3,b=4,c=5,k=5/12$；图中 $t_A<t_B<0$};
\begin{scope}[shift={(6.1,7)},scale=.8]
\coordinate (A) at (-4.1428571429,-3.8095238095);
\coordinate (B) at (3.0606060606,-.8080808081);
\coordinate (F) at (5,0);
\coordinate (H) at (3.0606060606,-3.8095238095);
\draw[gray,->] (-5.4,0)--(6,0) node[right] {$x$};
\draw[blue!75!black,thick,domain=-4.5:1,samples=90,variable=\v] plot ({3*sqrt(1+\v*\v/16)},\v);
\draw[blue!75!black,thick,domain=-4.5:1,samples=90,variable=\v] plot ({-3*sqrt(1+\v*\v/16)},\v);
\draw[red!75!black,dashed] (B)--(F);
\draw[red!75!black,very thick] (A)--(B) node[midway,above left=4pt] {$L$};
\draw[orange!85!black,thick] (A)--(H) node[midway,below=6pt] {$|t_A-t_B|$};
\draw[purple!80!black,thick] (H)--(B);
\node[anchor=west,text=purple!80!black] at (4.8,-2.3088023088) {$|k|\,|t_A-t_B|$};
\draw[gray] (2.8606060606,-3.8095238095)--(2.8606060606,-3.6095238095)--(3.0606060606,-3.6095238095);
\fill (A) circle(2pt) node[left=10pt] {$A$};
\fill (B) circle(2pt) node[left=10pt] {$B$};
\fill (F) circle(2pt) node[above right=5pt] {$F_2$};
\fill (H) circle(2pt) node[below right=5pt] {$H$};
\end{scope}
\node at (7.1,2.25) {$\displaystyle L^2=|t_A-t_B|^2+k^2|t_A-t_B|^2$};
\node at (7.1,1.2) {$\displaystyle L=\sqrt{1+k^2}\,|t_A-t_B|$};
\node at (7.1,.5) {$t=x-c$ 是横向增量；两根都为负，只说明两个交点都在焦点左侧。};
\end{tikzpicture}
\end{center}

过 $F_2$ 的非竖直直线 $y=k(x-c)$，令 $x=c+t$。
完整展开得到
\[
\frac{(c+t)^2}{a^2}-\frac{k^2t^2}{b^2}=1
\iff (b^2-a^2k^2)t^2+2cb^2t+b^4=0.
\]
利用 $c^2-a^2=b^2$，可写
\[
A_t=b^2-a^2k^2,\quad B_t=2cb^2,\quad C_t=b^4.
\]
由二次方程根与系数关系
\[
(t_1-t_2)^2=(t_1+t_2)^2-4t_1t_2
=\frac{B_t^2-4A_tC_t}{A_t^2}
=\frac{4a^2b^4(1+k^2)}{(b^2-a^2k^2)^2}.
\]
因为 $A,B$ 在同一条斜率为 $k$ 的直线上，
$|AB|=\sqrt{1+k^2}\,|t_1-t_2|$，得到结论一。

\tcbbreak
\textbf{方法三续：角度统一，先看斜率式漏掉的竖直弦。}
\begin{center}% Exact vertical focal chord, a=3,b=4,c=5; isotropic scale .5.
\begin{tikzpicture}[font=\small,>=stealth]
\draw[gray!45,rounded corners=2pt] (0,0) rectangle (14.2,10.3);
\node at (7.1,9.65) {角度法先补竖直：横坐标相同，仍然有两个不同交点};
\begin{scope}[shift={(6.2,5.2)},scale=.5]
\coordinate (A) at (5,5.3333333333);
\coordinate (B) at (5,-5.3333333333);
\coordinate (F) at (5,0);
\draw[blue!75!black,thick,domain=-6.2:6.2,samples=100,variable=\v] plot ({3*sqrt(1+\v*\v/16)},\v);
\draw[red!75!black,very thick] (B)--(A);
\draw[gray,->] (5,0)--(8,0);
\draw[orange!85!black,->] (5,0)--(5,2);
\draw[gray] (5.7,0) arc[start angle=0,end angle=90,radius=.7];
\node at (7,1.4) {$\theta=\pi/2$};
\fill (A) circle(2pt) node[left=8pt] {$A$};
\fill (B) circle(2pt) node[left=8pt] {$B$};
\fill (F) circle(2pt) node[below right=7pt] {$F_2$};
\end{scope}
\node[anchor=west] at (.8,7.3) {$x_A=x_B=c$};
\node[anchor=west] at (.8,6.2) {$\rho_A=b^2/a$};
\node[anchor=west] at (.8,5.1) {$\rho_B=-b^2/a$};
\node[anchor=west] at (.8,3.9) {$P=F_2+\rho(\cos\theta,\sin\theta)$};
\node at (7.1,1.4) {$\displaystyle |\rho_A-\rho_B|=\frac{2b^2}{a}=\frac{32}{3}$};
\node at (7.1,.5) {$\rho$ 为有向距离；$k$ 不存在，但倾斜角 $\theta=\pi/2$ 完全合法。};
\end{tikzpicture}
\end{center}
图中的 $\rho$ 是沿直线方向的有向距离，与 $t=x-c$ 不同；
竖直时 $t_A=t_B=0$，但两端点的 $\rho$ 不同，所以弦长仍为正。
\textbf{角度统一（补上竖直弦）：}先对 $\cos\theta\ne0$，令 $k=\tan\theta$，
则
\[
\begin{aligned}
L&=\frac{2ab^2(1+\tan^2\theta)}{|a^2\tan^2\theta-b^2|}
=\frac{2ab^2}{|a^2\sin^2\theta-b^2\cos^2\theta|}\\
&=\boxed{\frac{2ab^2}{|a^2-c^2\cos^2\theta|}}.
\end{aligned}
\]
当 $\cos\theta=0$ 时，\textbf{另用} $x=c$ 直接验算 $L=2b^2/a$；
因此角度表达式确实适用于 $\theta=\pi/2$。

\textbf{分支判别（给出严格的代数依据）：}
根积 $t_1t_2=b^4/(b^2-a^2k^2)$。
当 $|k|<b/a$，两根均负（根和也负），两交点位于左右不同支；
当 $|k|>b/a$，两根异号，焦点位于弦内，两交点在右支。
仅凭根的正负还不能直接判断左右支；严格的坐标范围由方法四末尾的显式根式验证。
\end{proofbox}

\clearpage
\begin{proofbox}
\textbf{方法四：补充｜焦半径＋余弦定理，解释同支最短为何在竖直方向}
\begin{center}% a=3,b=4,c=5,q=1/5; r=8,P=(33/5,16sqrt(6)/5),PF1=14.
\begin{tikzpicture}[font=\small,>=stealth]
\draw[gray!45,rounded corners=2pt] (0,0) rectangle (14.2,9.8);
\node at (7.1,9.15) {焦半径从哪来：先看底边 $2c$、两腰 $r$ 与 $r+2a$};
\node at (7.1,8.35) {示例 $a=3,b=4,c=5,\ \cos\varphi=1/5$};
\begin{scope}[shift={(5,2.6)},scale=.6]
\coordinate (Fone) at (-5,0);
\coordinate (Ftwo) at (5,0);
\coordinate (P) at (6.6,7.8383671769);
\draw[gray] (-5,0)--(8.3,0);
\draw[orange!85!black,very thick] (Fone)--(P) node[midway,above left=5pt] {$r+2a=14$};
\draw[red!75!black,very thick] (Ftwo)--(P) node[midway,right=6pt] {$r=8$};
\draw[blue!75!black,thick] (Fone)--(Ftwo) node[midway,below=6pt] {$2c=10$};
\draw[gray] (5.7,0) arc[start angle=0,end angle=78.4630409672,radius=.7];
\node at (6.1,.8) {$\varphi$};
\fill (Fone) circle(2pt) node[below left=6pt] {$F_1$};
\fill (Ftwo) circle(2pt) node[below right=6pt] {$F_2$};
\fill (P) circle(2pt) node[above left=6pt] {$P$};
\end{scope}
\node at (7.1,1.1) {$\angle F_1F_2P=\pi-\varphi$，故余弦定理中的交叉项为 $+4cr\cos\varphi$。};
\node at (7.1,.45) {$P$ 在右支：$PF_1-PF_2=2a$；$\varphi$ 从向右的水平射线开始量。};
\end{tikzpicture}
\end{center}

\textbf{增补理由：}不用背弦长公式，直接理解为何同支焦点弦的长度有 $2b^2/a$ 这个下界。
从右焦点 $F_2(c,0)$ 出发，沿与 $x$ 轴正向夹角 $\varphi$ 的射线取右支点 $P$，
记 $r=PF_2$。双曲线定义为 $PF_1=r+2a$；
三角形 $F_1F_2P$ 的内角 $\angle F_1F_2P=\pi-\varphi$，所以余弦定理中的交叉项为
$-4cr\cos(\pi-\varphi)=+4cr\cos\varphi$，得到
\[
PF_1^2=r^2+4cr\cos\varphi+4c^2.
\]
两式相等、消掉 $r^2$ 并用 $c^2-a^2=b^2$，得到
\[
\boxed{r=\frac{b^2}{a-c\cos\varphi}}.\tag{3}
\]
（这条式子仅在右支且分母为正时使用；不能无条件给左支套用。）

\tcbbreak
\textbf{方法四续：同支弦的两条反向射线，距离应当相加。}
\begin{center}% a=3,b=4,c=5,q=1/5; r+=8,r-=4; both endpoints on right branch.
\begin{tikzpicture}[font=\small,>=stealth]
\draw[gray!45,rounded corners=2pt] (0,0) rectangle (14.2,11.4);
\node at (7.1,10.75) {同支弦先看两条反向射线：余弦一正一负，距离相加};
\node at (7.1,9.85) {例 $a=3,b=4,c=5,q=1/5$：$r_+=8,\ r_-=4$};
\begin{scope}[shift={(4,4.1)},scale=.6]
\coordinate (A) at (6.6,7.8383671769);
\coordinate (B) at (4.2,-3.9191835885);
\coordinate (F) at (5,0);
\draw[blue!75!black,thick,domain=-4.5:8.3,samples=100,variable=\v] plot ({3*sqrt(1+\v*\v/16)},\v);
\draw[red!75!black,very thick] (B)--(A);
\draw[orange!85!black,->,thick] (F)--(5.5,2.4494897428);
\draw[orange!85!black,->,thick] (F)--(4.6,-1.9595917942);
\node[right=8pt] at (5.8,3.9191835885) {$r_+=8$};
\node[right=8pt] at (4.6,-1.9595917942) {$r_-=4$};
\fill (A) circle(2pt) node[left=8pt] {$A$};
\fill (B) circle(2pt) node[left=8pt] {$B$};
\fill (F) circle(2pt) node[below right=6pt] {$F_2$};
\end{scope}
\node[anchor=west] at (.8,7.3) {正向：$\cos\varphi=q$};
\node[anchor=west] at (.8,6.35) {反向：$\cos(\varphi+\pi)=-q$};
\node[anchor=west] at (.8,4.95) {$\displaystyle r_+=\frac{b^2}{a-cq}$};
\node[anchor=west] at (.8,3.55) {$\displaystyle r_-=\frac{b^2}{a+cq}$};
\node at (7.1,.7) {焦点在 $AB$ 内，所以 $L=r_++r_-$；当 $q=0$，两段相等、弦变竖直。};
\end{tikzpicture}
\end{center}
如果焦点 $F_2$ 被同支弦的两端点夹在中间，
它们沿反向射线，余弦分别是 $q$ 与 $-q$；故
\[
\boxed{L=\frac{b^2}{a-cq}+\frac{b^2}{a+cq}}
=\frac{2ab^2}{a^2-c^2q^2}\ge\frac{2b^2}{a},
\qquad |cq|<a.
\]
等号当且仅当 $q=\cos\varphi=0$，就是两端点都在\textbf{过焦点的竖直线}上。
这一步还解释了为何同支弦的分母没有负号：此时 $a^2-c^2q^2>0$。

\textbf{等价的根式补充（用于验算分支）：}令 $s=\sqrt{1+k^2}>0$，则二次式的根可分解为
\[
t_1=-\frac{b^2}{c+as},\qquad t_2=-\frac{b^2}{c-as},
\qquad(s\ne c/a).
\]
当 $1\le s<c/a$，$t_2\le-(c+a)$，故 $x_2=c+t_2\le-a$，对应左支；
$s=1$（即 $k=0$）恰好是左顶点。若 $s>c/a$，$t_2>0$，故 $x_2>c>a$，对应右支。
另有 $t_1\in[-(c-a),0)$，始终对应右支，故分类严格成立。
\end{proofbox}
